% Section 3: system architecture and data flow.
\section{System architecture}\label{sec:architecture}

\subsection{Layers}
\Cref{fig:arch} shows the package in four layers. The input layer turns MPS/QPS files, the
case-study builders or Python arrays into a \code{Model}. Presolve reduces the model and
scaling equilibrates it. Five engines solve it. All engines share one layer of sparse kernels, and
nothing below the API calls code outside the repository except NumPy, Numba and (for the GPU)
CuPy. \code{nirnay.solve(model, method="auto")} picks branch-and-bound for a MILP, the QP
interior-point method for a QP and the dual simplex for an LP. The interior-point method and
PDLP are chosen by name.

\begin{figure}[tbp]
\centering
\begin{tikzpicture}[x=1cm,y=1cm]
  % layer bands
  \foreach \y/\h/\name in {0/1.35/Input, -1.6/1.35/Presolve, -3.2/2.85/Engines, -6.3/1.6/Kernels}{
    \draw[layer] (0,\y) rectangle (15.4,\y-\h);
    \node[rotate=90,font=\sffamily\scriptsize\bfseries,text=navy] at (0.32,\y-\h/2) {\name};
  }
  % input
  \node[blk,minimum width=3.9cm] (mps) at (2.95,-0.675) {MPS / QPS reader\\[-1pt]{\scriptsize\color{muted}fixed, free, gzip, RANGES, QUADOBJ}};
  \node[blk,minimum width=4.2cm] (cases) at (7.45,-0.675) {Case-study builders\\[-1pt]{\scriptsize\color{muted}refinery, UC, dispatch, logistics}};
  \node[blk,minimum width=4.2cm] (api) at (12.4,-0.675) {Python API and CLI\\[-1pt]{\scriptsize\color{muted}\code{solve(model, method)}}};
  % presolve
  \node[blks,minimum width=6.6cm] (pre) at (4.4,-2.275) {Presolve and postsolve\\[-1pt]{\scriptsize singleton rows, fixed/empty columns, free singletons, duals}};
  \node[blks,minimum width=6.8cm] (scal) at (11.5,-2.275) {Scaling $\hat A = RAC$\\[-1pt]{\scriptsize geometric mean + equilibration, powers of two}};
  % engines
  \node[blkn,minimum width=2.55cm,minimum height=1.25cm] (spx) at (2.1,-4.25) {Dual simplex\\[-1pt]{\scriptsize DSE, BFRT, Harris}\\[-2pt]{\scriptsize phase 1, cleanup}};
  \node[blkn,minimum width=2.55cm,minimum height=1.25cm] (ipm) at (4.85,-4.25) {Interior point\\[-1pt]{\scriptsize Mehrotra PC}\\[-2pt]{\scriptsize normal equations}};
  \node[blkn,minimum width=2.55cm,minimum height=1.25cm] (pdlp) at (7.6,-4.25) {PDLP\\[-1pt]{\scriptsize restarted PDHG}\\[-2pt]{\scriptsize CPU \textbar{} GPU}};
  \node[blk,draw=igreen,fill=igreen,text=white,minimum width=2.9cm,minimum height=1.25cm] (bnb) at (11.0,-4.25) {Branch-and-bound\\[-1pt]{\scriptsize reliability branching}\\[-2pt]{\scriptsize propagation, GMI cuts, diving}};
  \node[blk,draw=plum,fill=plum,text=white,minimum width=2.45cm,minimum height=1.25cm] (qp) at (14.0,-4.25) {QP interior\\[-1pt] point\\[-2pt]{\scriptsize quasi-definite KKT}};
  \draw[->,igreen,line width=0.8pt,dashed] (bnb.north) -- ++(0,0.2) -| node[pos=0.25,above,lbl,text=igreen,yshift=-1pt]{node LPs: warm-started dual simplex} (spx.north);
  % kernels, each under the engine that uses it
  \foreach \k/\x/\t/\s in {lu/1.95/Sparse LU/Gilbert--Peierls + eta,
                          nrm/4.35/Normal matrix/{$A\Theta A\T$, fixed pattern},
                          chol/6.75/Cholesky/up-looking + etree,
                          spmv/9.15/SpMV kernels/Numba \textbar{} CUDA,
                          amd/11.55/Min. degree/mass elimination,
                          ldl/13.95/$LDL\T$/quasi-definite}{
    \node[blkm,minimum width=2.3cm,minimum height=1.0cm] (\k) at (\x,-7.0) {\t\\[-1pt]{\scriptsize\color{muted}\s}};
  }
  % engine -> kernel
  \draw[flow] ([xshift=-1.5mm]spx.south) -- ([xshift=-1.5mm]spx.south |- lu.north);
  \draw[flow] ([xshift=-3mm]ipm.south) -- ([xshift=-3mm]ipm.south |- nrm.north);
  \draw[flow] ([xshift=4mm]ipm.south) -- ++(0,-0.55) -| (chol.north);
  \draw[flow] ([xshift=4mm]pdlp.south) -- ++(0,-0.75) -| (spmv.north);
  \draw[flow] (qp.south) -- (qp.south |- ldl.north);
  \draw[flow,muted] (amd.east) -- (ldl.west);
  \draw[flow,muted] (amd.south) -- ++(0,-0.18) -| (chol.south);
  % input -> presolve -> engines
  \draw[flow] (mps.south) -- (mps.south |- pre.north);
  \draw[flow] (cases.south) -- (cases.south |- pre.north);
  \draw[flow] (api.south) -- (api.south |- scal.north);
  \draw[flow] (pre.east) -- (scal.west);
\end{tikzpicture}
\caption[System architecture]{System architecture. Everything in the four bands is \nirnay{} code.
The dashed arrow shows that branch-and-bound solves its node LPs with the warm-started dual
simplex. Minimum degree orders both the Cholesky and the $LDL\T$ factorisations.}
\label{fig:arch}
\end{figure}

\begin{table}[tbp]
\centering\small
\caption[Solver modules]{Solver modules of \nirnay{} with their line counts (read from the
source files by \file{gen\_data.py}). The core solver is \locCore{} lines; the case-study builders
add \locCases{}.}
\label{tab:modules}
\begin{tabularx}{\linewidth}{@{}l r X@{}}
\toprule
Module (\file{nirnay/\dots}) & Lines & Role\\
\midrule
\rows{data/modules.tex}
\bottomrule
\end{tabularx}
\end{table}

\subsection{The path of a solve}
\Cref{fig:dataflow} follows one LP or MILP from file to answer. Each transformation on the way
in has an exact inverse on the way out, so the returned $x$, $y$ and $z$ satisfy the optimality
conditions \eqref{eq:kkt-lp} of the model in the file, not only those of the scaled, reduced
model the engine saw.

\begin{figure}[tbp]
\centering
\begin{tikzpicture}[node distance=5mm and 5mm]
  \node[term] (file) {MPS file};
  \node[blk,right=of file] (read) {\code{read\_mps}};
  \node[blk,right=of read] (model) {\code{Model}\\[-1pt]{\scriptsize$c,A,r^l,r^u,\ell,u,\mathcal I,Q$}};
  \node[blks,right=of model] (pre) {presolve\\[-1pt]{\scriptsize stack of reductions}};
  \node[blks,right=of pre] (scale) {scale\\[-1pt]{\scriptsize$R,C=2^k$}};
  \node[blkn,below=11mm of scale,minimum width=2.6cm] (eng) {engine\\[-1pt]{\scriptsize simplex \textbar{} IPM \textbar{} PDLP \textbar{} B\&B}};
  \node[blks,left=of eng] (unscale) {unscale\\[-1pt]{\scriptsize$x=C\hat x,\ y=R\hat y$}};
  \node[blks,left=of unscale] (post) {postsolve\\[-1pt]{\scriptsize undo in reverse}};
  \node[term,fill=igreen,draw=igreen,left=of post] (res) {\code{Result}: $x,y,z$};
  \draw[flow] (file) -- (read);
  \draw[flow] (read) -- (model);
  \draw[flow] (model) -- (pre);
  \draw[flow] (pre) -- (scale);
  \draw[flow] (scale) -- (eng);
  \draw[flow] (eng) -- (unscale);
  \draw[flow] (unscale) -- (post);
  \draw[flow] (post) -- (res);
  % verification branch
  \node[blkg,below=9mm of res] (cmp) {compare: status, $\frac{|f-f^\star|}{\max(1,|f^\star|)}$,\\[-1pt]{\scriptsize row/bound violation, dual feasibility}};
  \node[blkm,right=8mm of cmp] (ref) {reference solver in its own process\\[-1pt]{\scriptsize HiGHS (LP, MILP) \textbar{} Clarabel (QP), own parser}};
  \draw[flow,muted] (file.west) -- ++(-0.3,0) |- ([yshift=-1.5mm]ref.south) -- (ref.south);
  \draw[flow,muted] (res.south) -- (cmp.north);
  \draw[flow,muted] (ref) -- (cmp);
  \begin{scope}[on background layer]
    \node[fit=(pre)(scale)(eng)(unscale)(post),draw=saffron!60,dashed,rounded corners=3pt,inner sep=5pt,
          label={[lbl,anchor=south east]north east:inside \code{nirnay.solve}}] {};
  \end{scope}
\end{tikzpicture}
\caption[Data flow of a solve]{Data flow of a solve. The upper row reduces and scales the model.
The engine solves it, and the lower row maps the solution back. The benchmark harness (bottom)
re-solves the same file with a reference solver in a separate process and checks the answer.}
\label{fig:dataflow}
\end{figure}

\subsection{Input: the MPS/QPS reader}
\file{nirnay/io/mps.py} reads fixed and free MPS, gzip-compressed files, \code{OBJSENSE},
integer \code{MARKER} blocks, \code{RANGES}, the bound types \code{UP LO FX FR MI PL BV LI UI},
an objective constant given as the RHS of the objective row, and the QP sections \code{QUADOBJ}
(one triangle), \code{QMATRIX} and \code{QSECTION} (both triangles). It follows the conventions the
benchmark libraries assume: an integer column in a \code{MARKER} block without bounds gets
$[0,\infty)$; \code{UP} with a negative value on a column whose lower bound is still $0$ sets
the lower bound to $-\infty$; \code{RANGES} on an equality row extend upward for $R>0$ and
downward for $R<0$; the objective constant is minus the RHS of the objective row. Fixed-format
files whose names contain spaces (Netlib's \code{forplan}) are detected and parsed by column
position. The earliest interior-point sweep failed to read \code{forplan} and \code{standgub}
(\cref{tab:fixes-ipm}); both read correctly now. A writer (\file{io/mps\_write.py}) emits
free MPS. Read, write and re-read reproduce the arrays of the benchmark files, and HiGHS reads the
original and the written file into the same model (\cref{sec:verification}).

\subsection{Scaling}\label{sec:scaling}
Industrial models mix units (tonnes, kilolitres, rupees, percentages), so the coefficients of
$A$ span many orders of magnitude and every factorisation loses digits to that spread.
\file{presolve/scaling.py} computes diagonal $R$ and $C$ and solves the scaled problem with
$\hat A=RAC$, $\hat c=Cc$, $\hat r = R r$ and $\hat\ell = C^{-1}\ell$, $\hat u = C^{-1}u$, in two
stages~\cite{elble2012}:
\begin{enumerate}
  \item \textbf{Geometric-mean passes.} Divide each row, then each column, by
        $\sqrt{\max_j|\hat a_{ij}|\cdot\min_j|\hat a_{ij}|}$. Stop after 8 passes, or earlier once
        a pass reduces the overall spread $\max|\hat a|/\min|\hat a|$ by less than 10\%. Entries
        below $10^{-12}$ of the largest in their row or column are left out of the minimum: they
        carry no information in double precision, but a geometric mean would let them dominate.
        The QP \code{KSIP} has an entry of order $10^{-30}$ next to entries of order~1, and scaling by
        it ruined every other entry of its line.
  \item \textbf{Equilibration.} A final column pass makes the largest entry of every column
        exactly~1.
\end{enumerate}
Each factor is then rounded to the nearest power of two, $R_i \leftarrow 2^{\operatorname{round}(\log_2 R_i)}$.
Multiplying by a power of two changes only the floating-point exponent, so scaling adds no rounding
error of its own, and unscaling ($x=C\hat x$, $y=R\hat y$, $z=C^{-1}\hat z$) returns exactly the
numbers the engine computed. PDLP uses its own preconditioner (Ruiz and Pock--Chambolle,
\cref{sec:pdlp}), because the convergence rate of a first-order method depends on the conditioning
of $A$ in a different way.
